\begin{abstract}
Les \q{ponts} d'O.~Caramello lisent une même propriété d'un topos sur deux présentations de ce topos, et tirent de ces deux lectures un théorème qui relie les présentations. Un article associé \cite{bridges-generator} fait de ces ponts un générateur de théorèmes, exécuté à la main. La présente note applique ce générateur aux notes manuscrites de Grothendieck (Archives Grothendieck, Université de Montpellier). Trois arches (deux valeurs, connexe, topos de préfaisceaux d'un préordre), démontrées en Lean~4 dans le paquet \texttt{lean/bridges} de cette édition, sont lues sur trois tabliers qui reviennent dans ses notes: la catégorie simpliciale, les groupes et les préordres. Nous discutons deux précédents tirés de ces notes: d'une part un exercice du dossier 161-3, p.~26, que le générateur retrouve par l'arche de Morita et dont la forme relative aux préordres est démontrée en Lean, avec la même identification au dossier 37, p.~48; d'autre part la correspondance entre sous-topos et quotients d'une théorie, au dossier 161-2, p.~17 et 44. Un dernier tablier, tiré du dossier 161-2, le topos des pseudo-corps (anneaux réguliers au sens de von Neumann), fait l'objet d'une section: nous vérifions l'une des trois présentations qu'il en donne, nous y lisons les arches et nous répondons à une question qu'il a laissée en marge. Une autre section reprend son lemme de recollement de sous-topos (dossier 161-2, p.~48) : nous l'énonçons et le démontrons, et montrons que dans le cas où il l'applique, les anneaux de caractéristique définie, il est en défaut, de sorte qu'aucun sous-topos du topos des anneaux ne classifie ces anneaux. Une annexe recense les fichiers Lean. Une version anglaise accompagne ce texte.
\end{abstract}
